\section{Welfare Sign Certificates}
\label{app:welfare}

This appendix supplies analytical sign certificates for the canonical welfare comparisons used in Section~\ref{sec:coalition-stability}. They are frozen relationships in the Stage-8 model; the calculations below make their domain-wide sign support explicit rather than relying on numerical sampling.

Recall
\[
W^{IS}=K_I+3P,
\qquad
W_M^{SU,N}=K_M+2A+C,
\qquad
W_O^{SU,N}=K_O+2B+D,
\qquad
W^{SW}=K_W+H+2S.
\]
The corrected international-standardization expression is
\[
W^{IS}=\frac{3(5-2v)}{32(1-v)^2}.
\]
All denominators below are strictly positive on the canonical domain.

\subsection{An SU member versus separate national standards}

Define the scaled welfare gap
\begin{align}
N_{MS}(c,v)
&\equiv 32(2-3v)^2\left(W_M^{SU,N}-W^{SW}\right)\notag\\
&=-216c^2v^2+264c^2v-60c^2
+108cv^2-144cv+56c
-99v^2+84v.
\label{eq:appendix-n-ms}
\end{align}
Its curvature in $c$ is
\[
\frac{\partial^2N_{MS}}{\partial c^2}
=-24(18v^2-22v+5)<0,
\]
because $18v^2-22v+5$ is decreasing on $0<v<1/4$ and equals $5/8$ at $v=1/4$. Hence the minimum over the closure of the admissible $c$ interval is attained at an endpoint. At $c=0$,
\[
N_{MS}(0,v)=3v(28-33v)>0.
\]
At $c_{\max}(v)=(1-3v)/[3(1-v)]$,
\[
N_{MS}(c_{\max},v)
=-\frac{p_{MS}(v)}{3(1-v)^2},
\]
where
\[
p_{MS}(v)=621v^4-1206v^3+729v^2-92v-36.
\]
For $0\le v\le1/4$, $v^4\le v^2/16$ and the term $-1206v^3$ is nonpositive, so
\[
p_{MS}(v)
\le \frac{12285}{16}v^2-92v-36.
\]
The quadratic on the right is convex and is negative at both endpoints of $[0,1/4]$; at $v=1/4$ it equals $-2819/256$. Therefore $p_{MS}(v)<0$, and
\begin{equation}
\boxed{W_M^{SU,N}>W^{SW}.}
\label{eq:appendix-su-member-gt-sw}
\end{equation}

\subsection{International standardization versus separate standards}

Similarly, let
\begin{align}
N_{IS}(c,v)
&\equiv 32(1-v)^2\left(W^{IS}-W^{SW}\right)\notag\\
&=-28c^2v^2+56c^2v-28c^2
+20cv^2-40cv+20c
-15v^2+24v.
\label{eq:appendix-n-is}
\end{align}
Since
\[
\frac{\partial^2N_{IS}}{\partial c^2}
=-56(1-v)^2<0,
\]
its minimum in $c$ is again at an endpoint. The endpoint values are
\[
N_{IS}(0,v)=3v(8-5v)>0
\]
and
\[
N_{IS}(c_{\max},v)
=\frac{32+144v-207v^2}{9}>0.
\]
The final inequality follows immediately from $v<1/4$, since $207v^2<207/16<32$ while $144v>0$. Thus
\begin{equation}
\boxed{W^{IS}>W^{SW}.}
\label{eq:appendix-is-gt-sw}
\end{equation}

\subsection{International standardization versus the SU outsider}

Finally define
\begin{align}
N_{IO}(c,v)
&\equiv 32(1-v)^2(2-3v)^2
\left(W^{IS}-W_O^{SU,N}\right).
\label{eq:appendix-n-io-definition}
\end{align}
The exact polynomial is
\begin{align}
N_{IO}(c,v)
={}&-144c^2v^4+576c^2v^3-912c^2v^2+672c^2v-192c^2\notag\\
&+288cv^4-912cv^3+1072cv^2-560cv+112c\notag\\
&-252v^4+666v^3-537v^2+132v.
\label{eq:appendix-n-io}
\end{align}
Its curvature is
\[
\frac{\partial^2N_{IO}}{\partial c^2}
=-96(1-v)^2(3v^2-6v+4)<0.
\]
At $c=0$,
\[
N_{IO}(0,v)
=-3v\,p_0(v),
\qquad
p_0(v)=84v^3-222v^2+179v-44.
\]
The derivative $p_0'(v)=252v^2-444v+179$ is decreasing on $[0,1/4]$ and remains positive at $v=1/4$, where it equals $335/4$. Hence $p_0$ is increasing, while $p_0(1/4)=-189/16<0$; therefore $N_{IO}(0,v)>0$.

At the upper $c$ boundary,
\[
N_{IO}(c_{\max},v)
=-\frac{p_1(v)}{3},
\]
where
\[
p_1(v)=324v^4-990v^3+843v^2-92v-48.
\]
Using $v^4\le v^2/16$ and dropping the nonpositive cubic term gives
\[
p_1(v)
\le \frac{3453}{4}v^2-92v-48.
\]
This bounding quadratic is convex and negative at both endpoints of $[0,1/4]$; at $v=1/4$ it equals $-1091/64$. Thus $p_1(v)<0$ and the upper-boundary value of $N_{IO}$ is positive. Concavity in $c$ now yields
\begin{equation}
\boxed{W^{IS}>W_O^{SU,N}.}
\label{eq:appendix-is-gt-su-outsider}
\end{equation}

\subsection{The outsider post-bypass gap}
\label{app:j-positive}

Section~\ref{sec:coalition-stability} defines
\[
J=K_I+P-K_O-D,
\]
so that $W^{IS}-W_O^{SU,O}=J+F$. To establish the strict outsider ranking after bypass it is sufficient to show $J>0$. Multiplying $J$ by its positive denominator gives the sign-equivalent numerator
\begin{align}
N_J(c,v)
={}&-48(1-v)^2c^2
+(48v^3-80v^2+16v+16)c\notag\\
&+v(68-269v+318v^2-108v^3).
\label{eq:appendix-j-numerator}
\end{align}
For fixed $v$, $N_J$ is concave in $c$ because
\[
\frac{\partial^2N_J}{\partial c^2}=-96(1-v)^2<0.
\]
Its minimum on the closure of the admissible $c$ interval is therefore attained at an endpoint. At $c=0$ and $c=c_{\max}(v)$,
\begin{align}
N_J(0,v)
&=v(68-269v+318v^2-108v^3),\notag\\
N_J(c_{\max},v)
&=\frac v3(284-1095v+1098v^2-324v^3).
\label{eq:appendix-j-endpoints}
\end{align}
Both bracketed polynomials are decreasing on $0<v<1/4$ and remain positive at $v=1/4$, where they equal $303/16$ and $1181/16$. Hence $N_J>0$ and therefore
\begin{equation}
\boxed{J>0.}
\label{eq:appendix-j-positive}
\end{equation}
Since $F>0$, this implies $W^{IS}>W_O^{SU,O}$ throughout the canonical domain whenever the outsider-only continuation is relevant.

Equations~\eqref{eq:appendix-su-member-gt-sw}--\eqref{eq:appendix-j-positive} are precisely the sign comparisons invoked in the strict-blocking analysis. The symbolic verification scripts check the scaled-polynomial representations, curvature expressions, and endpoint identities independently from the manuscript text.
